\setcounter{exercisegroup}{8}
\edef\CurrentExerciseGroup{\arabic{exercisegroup}}
\section*{\texorpdfstring{\S\CurrentExerciseGroup}{Section \CurrentExerciseGroup}}
\phantomsection
\addcontentsline{toc}{section}{\texorpdfstring{Exercises for \S\CurrentExerciseGroup}{Exercises for Section \CurrentExerciseGroup}}

\begin{enumerate}
\def\labelenumi{\arabic{enumi})}
\item
  Let \(\mathrm{T}, \mathrm{T}'\) be two locally compact spaces, \(\mu\) a positive measure on \(\mathrm{T}\) , \(\mu'\) a positive measure on \(\mathrm{T}'\) . Show that if the measure \(\mu'\) is bounded, then the projection \(\mathrm{pr}_1\) of \(\mathrm{T} \times \mathrm{T}'\) onto \(\mathrm{T}\) is a \((\mu \otimes \mu')\) -proper mapping and \(\mathrm{pr}_1(\mu \otimes \mu') = a \cdot \mu\) , where \(a = \mu'(\mathrm{T}')\) . Deduce from this an example of a \((\mu \otimes \mu')\) -negligible subset of \(\mathrm{T} \times \mathrm{T}'\) whose projection on \(\mathrm{T}\) is not \(\mu\) -measurable.
\item
  Let T be the interval \([0,1]\) of R, equipped with the topology induced by that of R, and \(T'\) the interval \([0,1]\) of R, equipped with the discrete topology; let \(\mu\) be Lebesgue measure on T, \(\mu'\) the discrete measure on \(T'\) defined by the mass +1 at every point of \(T'\) , and let \(\nu = \mu \otimes \mu'\) be the product measure on \(X = T \times T'\) .
\end{enumerate}

\begin{enumerate}
\def\labelenumi{\alph{enumi})}
\item
  For every \(t \in \mathrm{T}\) , the set \(\{t\}\) is \(\mu\) -negligible, but the set \(\{t\} \times \mathrm{T}'\) is not \(\nu\) -negligible.
\item
  The `diagonal' \(\Delta\) of \(\mathrm{T} \times \mathrm{T}'\) (the set of points \((t, t)\) , where \(t\) runs over \([0, 1]\) ) is a closed set in \(\mathrm{T} \times \mathrm{T}'\) , locally \(\nu\) -negligible but not \(\nu\) -negligible, such that \(\mu'(\Delta(t)) = 1\) for all \(t \in \mathrm{T}\) and \(\mu(\Delta(t')) = 0\) for all \(t' \in \mathrm{T}'\) (cf.~§3, Exer. 3).
\item
  On the space \(\mathrm{X} \times \mathrm{X} = \mathrm{Y}\) , consider the product measure \(\rho = \nu \times \nu\) . Show that there exists in \(\mathrm{Y}\) a set \(\mathrm{A}\) , locally \(\rho\) -negligible, such that for every \(x \in \mathrm{X}\) , \(\mathrm{A}(x)\) and \(\mathrm{A}^{-1}(x)\) are \(\nu\) -measurable and not locally negligible. (If \(x = (t, t')\) , where \(t \in \mathrm{T}\) , \(t' \in \mathrm{T}'\) , take A to be such that A(x) contains all of the points \((s, t) \in X\) , where \(0 \leqslant s \leqslant 1\) .)
\end{enumerate}

\begin{enumerate}
\def\labelenumi{\arabic{enumi})}
\setcounter{enumi}{2}
\tightlist
\item
  Let T be a locally compact space, \(\mu\) a positive measure on T, f a numerical function \(\geqslant 0\) defined on T. In the product space \(T \times R\) , denote by \(D_{f}\) the set of points \((t, x)\) such that \(0 \leqslant x \leqslant f(t)\) ; on the other hand let \(\nu\) be Lebesgue measure on R.
\end{enumerate}

\begin{enumerate}
\def\labelenumi{\alph{enumi})}
\item
  Show that for \(f\) to be \(\mu\) -measurable, it is necessary and sufficient that \(D_f\) be a measurable set for the product measure \(\lambda = \mu \otimes \nu\) . (To see that the condition is necessary, prove that if \(f\) is \(\mu\) -measurable, then for every compact subset \(K\) of \(T\) , \(D_f \cap (K \times R_+)\) is the union of a \(\lambda\) -negligible set and a countable family of sets of the form \(A \times I\) , where \(A\) is \(\mu\) -measurable and \(I\) is an interval of \(R_+\) . To show that the condition is sufficient, prove that if it is satisfied then, for every compact subset \(K\) of \(T\) , there exists a dense set \(H \subset R_+\) such that, for every \(\alpha \in H\) , \(K \cap f^{-1}([\alpha, +\infty])\) is \(\mu\) -measurable, using Cor. 2 of Prop. 7.)
\item
  Show that, for \(f\) to be \(\mu\) -integrable, it is necessary and sufficient that \(\mathrm{D}_f\) be \(\lambda\) -integrable, and one then has \(\lambda(\mathrm{D}_f) = \int f d\mu\) . Moreover, denoting by \(g\) the decreasing numerical function on \(\mathbf{R}_+\) defined by \(g(t) = \mu\left( f([t, +\infty]) \right)\) (possibly equal to \(+\infty\) for \(t = 0\) (Ch. IV, §5, Exer. 29)), one has \(\int f d\mu = \int_0^{+\infty} g(t) dt\) .
\end{enumerate}

\begin{enumerate}
\def\labelenumi{\arabic{enumi})}
\setcounter{enumi}{3}
\tightlist
\item
  \begin{enumerate}
  \def\labelenumii{\alph{enumii})}
  \tightlist
  \item
    Let \(T, T'\) be two locally compact spaces countable at infinity, \(\mu\) a positive measure on T, \(\mu'\) a positive measure on \(T'\) . Show that for every \((\mu \otimes \mu')\) -measurable mapping f of \(T \times T'\) into R, the function \(F : t \mapsto \int^{*} |f(t, t')| d\mu'(t')\) is \(\mu\) -measurable (consider \(T'\) as the union of an increasing sequence of compact sets).
  \end{enumerate}
\end{enumerate}

\begin{enumerate}
\def\labelenumi{\alph{enumi})}
\setcounter{enumi}{1}
\tightlist
\item
  Suppose, moreover, that for almost every \(t \in T\) (for \(\mu\) ) the function \(f(t, \cdot)\) is \(\mu'\) -integrable and that for almost every \(t' \in T'\) (for \(\mu'\) ) the function \(f(\cdot, t')\) is \(\mu\) -integrable; finally, assume that the function (defined almost everywhere for \(\mu\) )
\end{enumerate}

\(t \mapsto \int f(t, t') d\mu'(t')\) is \(\mu\) -integrable. Show that the function (defined almost everywhere for \(\mu'\) ) \(t' \mapsto \int f(t, t') d\mu(t)\) is then \(\mu'\) -integrable and

\[
\int d \mu^ {\prime} (t ^ {\prime}) \int f (t, t ^ {\prime}) d \mu (t) = \int d \mu (t) \int f (t, t ^ {\prime}) d \mu^ {\prime} (t ^ {\prime}).
\]

(By virtue of \(a\) ), there is a partition of \(\mathbf{T}\) formed by a \(\mu\) -negligible set \(\mathbf{N}\) and a sequence \((\mathrm{K}_n)\) of compact sets such that each of the functions \(f\varphi_{\mathrm{K}_n\times \mathrm{T}'}\) is \((\mu \otimes \mu')\) -integrable. Set

\[
g _ {n} (t ^ {\prime}) = \sum_ {j = 1} ^ {n} \int_ {K _ {j}} f (t, t ^ {\prime}) d \mu (t)
\]

and apply Exer. 20 of §5.)

\begin{enumerate}
\def\labelenumi{\alph{enumi})}
\setcounter{enumi}{2}
\tightlist
\item
  Let \(\mathbf{T}\) be the interval \([0,1]\) of \(\mathbf{R}\) , \(\mu\) Lebesgue measure on \(\mathbf{T}\) . For every integer \(n > 0\) , let \(\mathrm{A}_n^\prime = \left[\frac{1}{2^n}, \frac{3}{2^{n+1}}\right]\) , \(\mathrm{A}_n^{\prime\prime} = \left[\frac{3}{2^{n+1}}, \frac{1}{2^{n-1}}\right]\) , and, in the product space \(\mathbf{T} \times \mathbf{T}\) , let \(\mathrm{B}_n^\prime = \mathrm{A}_n^\prime \times \mathrm{A}_n^\prime\) , \(\mathrm{B}_n^{\prime\prime} = \mathrm{A}_n^{\prime\prime} \times \mathrm{A}_n^{\prime\prime}\) , \(\mathrm{C}_n^\prime = \mathrm{A}_n^\prime \times \mathrm{A}_n^{\prime\prime}\) , \(\mathrm{C}_n^{\prime\prime} = \mathrm{A}_n^{\prime\prime} \times \mathrm{A}_n^\prime\) . Set \(f(t,t') = 4^{n+1}\) in \(\mathrm{B}_n^\prime\) and \(\mathrm{B}_n^{\prime\prime}\) , \(f(t,t') = -4^{n+1}\) in \(\mathrm{C}_n^\prime\) and \(\mathrm{C}_n^{\prime\prime}\) , for every integer \(n > 0\) , and \(f(t,t') = 0\) at the other points of \(\mathbf{T} \times \mathbf{T}\) . Show that the two integrals \(\int d\mu(t') \int f(t,t') d\mu(t)\) and \(\int d\mu(t) \int f(t,t') d\mu(t')\) are defined and equal, but the function \(f\) is not integrable for the measure \(\mu \otimes \mu\) .
\end{enumerate}

¶ 5) Let T, T' be two locally compact spaces, μ a positive measure on T, μ' a positive measure on T'; suppose that every compact subset of T is metrizable. Let f be a mapping of T × T' into a metrizable space G, such that: 1° for every t ∈ T, the mapping t' → f(t, t') is μ'-measurable; 2° for every t' ∈ T', the mapping t → f(t, t') is continuous. Show that, under these conditions, f is (μ ⊗ μ')-measurable. (Observe (using Egoroff's theorem) that for every compact subset K of T, the restriction of f to K × T' is the limit of a sequence of (μ ⊗ μ')-measurable functions.)

¶ 6) a) Let T be a locally compact space, ν a positive measure on T, I an interval of R, μ a positive measure on I. Let A be a subset of the product I × T, such that: 1° for every x ∈ I, the section A(x) ⊂ T of A at x is ν-measurable; 2° the relation x ≤ y implies A(x) ⊂ A(y). Show that A is (μ ⊗ ν)-measurable. (Reduce to the case that I and T are compact and μ is diffuse; consider an increasing sequence (xi)0≤i≤n of points of I, x0 and xn being the end-points of I, such that μ({[}xi, xi+1{]}) ≤ ε for 0 ≤ i ≤ n-1, and form two subsets B,C of I × T, measurable for λ = μ ⊗ ν and such that B ⊂ A ⊂ C and λ(C - B) ≤ εν(t).)

\begin{enumerate}
\def\labelenumi{\alph{enumi})}
\setcounter{enumi}{1}
\item
  Deduce from a) that if \(f\) is a numerical function defined on \(\mathrm{I} \times \mathrm{T}\) , such that: \(1^{\circ}\) for each \(x \in \mathrm{I}\) , \(t \mapsto f(x, t)\) is \(\nu\) -measurable, and \(2^{\circ}\) for each \(t \in \mathrm{T}\) , \(x \mapsto f(x, t)\) is increasing, then \(f\) is \((\mu \otimes \nu)\) -measurable.
\item
  Let \(g\) be a numerical function defined on \(I \times T\) , such that: \(1^{\circ}\) for each \(x \in I\) , \(t \mapsto g(x, t)\) is \(\nu\) -measurable; \(2^{\circ}\) for each \(t \in T\) , \(x \mapsto g(x, t)\) is monotone. Deduce from \(b\) ) that \(g\) is \((\mu \otimes \nu)\) -measurable. (Let \(T_1 \subset T\) be the set of \(t \in T\) such that the mapping \(x \mapsto g(x, t)\) is increasing; show that \(T_1\) is \(\nu\) -measurable. To that end, for every pair \((r_1, r_2)\) of rational numbers belonging to \(I\) such that \(r_1 \leqslant r_2\) , consider the set \(D_{r_1, r_2}\) of points \(t \in T\) such that \(g(r_1, t) \leqslant g(r_2, t)\) and express \(T_1\) as the intersection of sets \(D_{r_1, r_2}\) .)
\end{enumerate}

¶ 7) a) Let T, T' be two locally compact spaces; assume that T' admits a countable base. Let μ be a positive measure on T, μ' a positive measure on T'. Let f be a numerical function ≥0 defined on T × T', bounded on every compact subset of T × T', and such that: \(1^{\circ}\) for almost every \(t \in \mathrm{T}\) , the function \(t' \mapsto f(t, t')\) is \(\mu'\) -measurable; \(2^{\circ}\) for every function \(h \in \mathcal{K}(\mathrm{T}')\) , the function

\[
t \mapsto \int f (t, t ^ {\prime}) h (t ^ {\prime}) d \mu^ {\prime} (t ^ {\prime}),
\]

defined almost everywhere, is \(\mu\) -measurable. Show that, under these conditions, there exists a function \(g\) , measurable for \(\mu \otimes \mu'\) , such that for every \(t \in \mathrm{T}\) , one has \(f(t, t') = g(t, t')\) except at the points of a \(\mu'\) -negligible set \(\mathrm{A}_t\) (depending on \(t\) ). (Show that for every function \(\varphi \in \mathcal{K}(\mathrm{T} \times \mathrm{T}')\) , the function \(t' \mapsto f(t, t')\varphi(t, t')\) is \(\mu'\) -integrable for almost every \(t \in \mathrm{T}\) , and that the function \(t \mapsto \int f(t, t')\varphi(t, t') d\mu'(t')\) , defined almost everywhere, is \(\mu\) -integrable; one makes use of Lemma 1 of Ch. III, §4, No.~1. Observe next that \(\varphi \mapsto \int d\mu(t) \int f(t, t')\varphi(t, t') d\mu'(t')\) is a positive measure on \(\mathrm{T} \times \mathrm{T}'\) , with base \(\mu \otimes \mu'\) , and apply the Lebesgue-Nikodym theorem; finally, use the fact that there exists a countable dense set in \(\mathcal{K}(\mathrm{T}')\) .)

\begin{enumerate}
\def\labelenumi{\alph{enumi})}
\setcounter{enumi}{1}
\item
  Suppose, moreover, that T is metrizable. Show that the conditions of a) are satisfied if: \(1^{\circ}\) for almost every \(t' \in T'\) , the function \(t \mapsto f(t, t')\) is \(\mu\) -measurable; \(2^{\circ}\) for almost every \(t \in T\) , the function \(t' \mapsto f(t, t')\) is continuous almost everywhere (for \(\mu'\) ). (Use Exer. 13 of Ch. IV, §5.)
\item
  Take \(\mathrm{T} = \mathrm{T}' = [0,1]\) . Admitting the continuum hypothesis (S, Ch. III, §6, No.~4), let \(x \prec y\) be a well-ordering relation on \(\mathrm{T}\) for which there is no greatest element, and such that for every \(x \in \mathrm{T}\) , the set of \(z \prec x\) is countable. If one takes for \(\mu = \mu'\) Lebesgue measure, show that the characteristic function \(f\) of the set of pairs \((t, t')\) such that \(t \prec t'\) satisfies the conditions of \(a\) ), but that \(\int f(t, t') d\mu(t) = 0\) for every \(t' \in \mathrm{T}'\) , and \(\int f(t, t') d\mu'(t') = 1\) for every \(t \in \mathrm{T}\) .
\end{enumerate}

\begin{enumerate}
\def\labelenumi{\arabic{enumi})}
\setcounter{enumi}{7}
\tightlist
\item
  Let \(\mathrm{T}\) be a locally compact space, \(\mu\) a positive measure \(\neq 0\) on \(\mathrm{T}\) , A a discrete space, \(\lambda\) the measure on \(\mathrm{A}\) defined by mass \(+1\) at each point of \(\mathrm{A}\) .
\end{enumerate}

\begin{enumerate}
\def\labelenumi{\alph{enumi})}
\item
  For a mapping \(f\) of \(\mathrm{A} \times \mathrm{T}\) into a topological space to be \((\lambda \otimes \mu)\) -measurable, it is necessary and sufficient that, for every \(\alpha \in \mathrm{A}\) , the mapping \(t \mapsto f(\alpha, t)\) be \(\mu\) -measurable.
\item
  For a function \(\mathbf{f}\) defined on \(\mathrm{A} \times \mathrm{T}\) , with values in \(\overline{\mathbf{R}}\) or in a Banach space, to be \((\lambda \otimes \mu)\) -integrable, it is necessary and sufficient that: \(1^{\circ}\) except for a countable set of values of \(\alpha \in \mathrm{A}\) , the function \(t \mapsto \mathbf{f}(\alpha, t)\) be identically zero; \(2^{\circ}\) for every \(\alpha \in \mathrm{A}\) , the function \(t \mapsto \mathbf{f}(\alpha, t)\) be \(\mu\) -integrable; \(3^{\circ} \sum_{\alpha \in \mathrm{A}} \int |\mathbf{f}(\alpha, t)| d\mu(t) < +\infty\) . For \(\mathbf{f}\) to be
\end{enumerate}

essentially \((\lambda \otimes \mu)\) -integrable, it necessary and sufficient that: \(1^{\circ}\) for each \(\alpha \in \mathbf{A}\) , the function \(t \mapsto \mathbf{f}(\alpha, t)\) be essentially \(\mu\) -integrable; \(2^{\circ} \sum_{\alpha \in \mathbf{A}} \int |\mathbf{f}(\alpha, t)| d\mu(t) < +\infty\) ; one then

has

\[
\iint \mathbf {f} d \lambda d \mu = \sum_ {\alpha \in A} \int \mathbf {f} (\alpha , t) d \mu (t).
\]

\begin{enumerate}
\def\labelenumi{\alph{enumi})}
\setcounter{enumi}{2}
\tightlist
\item
  Let X be a locally compact space, \((\alpha,t)\mapsto\rho_{\alpha,t}\) a family of positive measures on X \((\alpha\in\mathrm{A},t\in\mathrm{T})\) . In order that the family \((\alpha,t)\mapsto\rho_{\alpha,t}\) be \((\mu\otimes\nu)\) -adequate, it is necessary and sufficient that, for every \(\alpha\in\mathrm{A}\) , the family \(t\mapsto\rho_{\alpha,t}\) be \(\mu\) -adequate, and that the family of measures \(\left(\int\rho_{\alpha,t}d\mu(t)\right)_{\alpha\in\mathrm{A}}\) be summable.
\end{enumerate}

\begin{enumerate}
\def\labelenumi{\arabic{enumi})}
\setcounter{enumi}{8}
\tightlist
\item
  Let \(u\) and \(v\) be two increasing and right-continuous numerical functions on \(\mathbf{R}\) , such that \(u(x) = v(x) = 0\) for \(x < 0\) . Let \(w\) be the increasing right-continuous function on \(\mathbf{R}\) defined by \(w(t) = u(t)v(t)\) for \(t \geqslant 0\) , \(w(t) = 0\) for \(t < 0\) ; let \(\lambda, \mu, \nu\) be the Stieltjes measures associated with \(u, v, w\) respectively (§6, Exer. 5).
\end{enumerate}

To every function \(\mathbf{f}\) defined on \(\mathbf{R}\) , with values in \(\overline{\mathbf{R}}\) or in a Banach space \(\mathrm{F}\) , one makes correspond the function \(\overline{\mathbf{f}}\) defined on \(\mathbf{R}^2\) by the conditions \(\overline{\mathbf{f}}(x,y) = \mathbf{f}(x)\) if \(y < x\) , \(\overline{\mathbf{f}}(x,y) = \mathbf{f}(y)\) if \(y \geqslant x\) . Show that, for \(\mathbf{f}\) to be \(\nu\) -integrable, it is necessary and sufficient that \(\overline{\mathbf{f}}\) be integrable for the product measure \(\lambda \otimes \mu\) , in which case \(\int \mathbf{f} d\nu = \iint \overline{\mathbf{f}} d\lambda d\mu\) (prove it first for the characteristic functions of intervals). From this, deduce the formula

\[
\int \mathbf {f} (x) d w (x) = \int \mathbf {f} (x) v (x -) d u (x) + \int \mathbf {f} (x) u (x +) d v (x).
\]

In particular, if \(u\) and \(v\) are continuous on \(\mathbf{R}\) , one obtains the formula for integration by parts

\[
\int_ {a} ^ {b} u (x) d v (x) = u (b) v (b) - u (a) v (a) - \int_ {a} ^ {b} v (x) d u (x).
\]

\begin{enumerate}
\def\labelenumi{\arabic{enumi})}
\setcounter{enumi}{9}
\tightlist
\item
  Let \(\mathbf{T}\) and \(\mathbf{X}\) be two locally compact spaces, \(\lambda\) a positive measure on \(\mathbf{T}\) , \(\mu\) a bounded positive measure on \(\mathbf{X}\) such that \(\mu(\mathbf{X}) = 1\) . Let \(f\) be a numerical function \(> 0\) at every point of \(\mathbf{T} \times \mathbf{X}\) , integrable along with \(\log f\) for the measure \(\lambda \otimes \mu\) . Prove the inequality
\end{enumerate}

\[
\log \left(\int \exp (\int \log f d \mu) d \lambda\right) \leqslant \int (\log \int f d \lambda) d \mu
\]

and show that equality cannot hold unless \(f\) is equivalent to a function of the form \(g \otimes h\) (apply the inequality of the geometric mean (Ch. IV, §6, Exer. 7 d)), for every \(t \in \mathrm{T}\) , to the function \(x \mapsto f(t,x) / (\int f(t,x)d\lambda (t))\) .

\begin{enumerate}
\def\labelenumi{\arabic{enumi})}
\setcounter{enumi}{10}
\tightlist
\item
  Let \(p\) be a finite real number \(\geqslant 1\) , \(T\) and \(X\) two locally compact spaces, \(\lambda\) a positive measure on \(T\) , \(\mu\) a positive measure on \(X\) , and \(f\) a function \(\geqslant 0\) defined on \(T \times X\) , integrable along with \(f^p\) for the measure \(\lambda \otimes \mu\) . Prove the inequality
\end{enumerate}

\[
\left(\int^ {*} \left(\int f d \mu\right) ^ {p} d \lambda\right) ^ {1 / p} \leqslant \int^ {*} \left(\int f ^ {p} d \lambda\right) ^ {1 / p} d \mu .
\]

(For every \(t \in \mathbf{T}\) , apply Hölder's inequality to the function \(x \mapsto f(t, x)\) , put into the form

\[
f (t, x) = g (t, x) \left(\int f ^ {p} (t, x) d \lambda (t)\right) ^ {1 / p q},
\]

where \(q\) is the exponent conjugate to \(p\) .) Show that equality cannot hold unless \(f\) is equivalent to a function of the form \(g \otimes h\) .

¶ 12) Let \(\mathrm{T}_i\) ( \(1 \leqslant i \leqslant n\) ) be \(n\) locally compact spaces, \(\mu_i\) a positive measure on \(\mathrm{T}_i\) for \(1 \leqslant i \leqslant n\) . For each index \(i\) , denote by \(\mathrm{E}_i\) the product \(\prod_{j \neq i} \mathrm{T}_j\) ; let \(f_i\)

be a function \(\geqslant 0\) on \(\mathrm{T} = \prod_{i=1}^{n} \mathrm{T}_i\) , measurable for \(\mu = \bigotimes_{i=1}^{n} \mu_i\) and not dependent

on \(t_{i}\) ; show that if, for \(1 \leqslant k \leqslant n\) , the function \(f_{k}^{n-1}\) is integrable for the measure \(\mu_{1} \otimes \cdots \otimes \mu_{k-1} \otimes \mu_{k+1} \otimes \cdots \otimes \mu_{n}\) , then \(f_{1}f_{2} \cdots f_{n}\) is \(\mu\) -integrable, and

\[
\int f _ {1} f _ {2} \dots f _ {n} d \mu_ {1} d \mu_ {2} \dots d \mu_ {n} \leqslant \left(\prod_ {k = 1} ^ {n} \mathrm{J} _ {k}\right) ^ {1 / (n - 1)}
\]

where, for each index k, one sets \(J_{k}=\int f_{k}^{n-1}d\mu_{1}\cdots d\mu_{k-1}d\mu_{k+1}\cdots d\mu_{n}\) (proceed by induction on n, applying the Lebesgue--Fubini theorem and Hölder's inequality).

Deduce from this that if A is a \(\mu\) -measurable subset of T, \(A_{i}\) its projection on \(E_{i}\) , and if \(A_{i}\) is integrable and of measure \(m_{i}\) (for the measure \(\bigotimes_{j\neq i}\mu_{j}\) on \(E_{i}\) ), then A is

\(\mu\) -integrable and

\[
\mu (\mathrm{A}) \leqslant (m _ {1} m _ {2} \dots m _ {n}) ^ {1 / (n - 1)}.
\]

Examine the case of equality in these two inequalities.

Generalize to the case where, instead of considering the \(n\) products of \(n - 1\) of the \(\mathrm{T}_i\) , one considers the \(\binom{n}{p}\) products of \(p\) of the \(\mathrm{T}_i\) , and where one integrates over \(\mathrm{T}\) a product of \(\binom{n}{p}\) functions \(\geqslant 0\) each of which depends on only \(p\) of the variables \(t_i\) . For example, if \(\mathrm{F}_{ij} = \mathrm{T}_i \times \mathrm{T}_j (i < j)\) and if \(f_{ij}\) depends only on the variables \(t_i\) and \(t_j\) , show that

\[
\int \left(\prod_ {i <   j} f _ {i j}\right) d \mu_ {1} d \mu_ {2} \dots d \mu_ {n} \leqslant \left(\prod_ {i <   j} \int f _ {i j} ^ {n - 1} d \mu_ {i} d \mu_ {j}\right) ^ {1 / (n - 1)}.
\]

¶ 13) Let \((T_{\iota})_{\iota \in I}\) be a family of compact spaces, and for each \(\iota \in I\) , let \(\mu_{\iota}\) be a positive measure on \(T_{\iota}\) of total mass equal to 1. Let \(T\) be the product space \(\prod_{\iota \in I} T_{\iota}\) ,

and \(\mu\) the product measure \(\bigotimes \mu_{\iota}\) on T (Ch. III, §4, No.~5).

\begin{enumerate}
\def\labelenumi{\alph{enumi})}
\item
  For every \(\iota \in I\) , let \(K_{\iota}\) be a compact subset of \(T_{\iota}\) ; show that if \(K = \prod_{\iota \in I} K_{\iota}\) then \(\mu(K) = \prod \mu_{\iota}(K_{\iota})\) .
\item
  Let M be a compact subset of T. Show that there exists a countable subset J of I such that, setting H = I - J, \(T_{J} = \prod_{\iota \in J} T_{\iota}\) , \(\mu_{J} = \bigotimes_{\iota \in J} \mu_{\iota}\) , \(T_{H} = \prod_{\iota \in H} T_{\iota}\) , one has \(M \subset N \times T_{H}\) , where N is a \(\mu_{J}\) -measurable subset of \(T_{J}\) such that \(\mu_{\mathrm{J}}(\overline{\mathrm{N}}) = \mu(\mathrm{M})\) . (Observe that for every n \textgreater{} 0, there exists an open neighborhood of M, the union of a finite number of elementary sets, whose measure differs from \(\mu(\mathrm{M})\) by less than 1/n.) From this, deduce that for almost every \(x \in N\) , the section \(\mathrm{M}(x) \subset \mathrm{T}_{\mathrm{H}}\) (when T is identified with the product \(T_{J} \times T_{H}\) ) contains the product of the supports of the measures \(\mu_{\iota}\) for \(\iota \in H\) (make use of the Lebesgue--Fubini theorem).
\item
  Show that if I is countable and if, for every \(\iota \in I\) , \(A_{\iota}\) is a \(\mu_{\iota}\) -measurable subset of \(T_{\iota}\) , then the set \(A = \prod_{\iota \in I} A_{\iota}\) is \(\mu\) -measurable and \(\mu(A) = \prod_{\iota \in I} \mu_{\iota}(A_{\iota})\) .
\item
  Suppose that I is uncountable. For every \(\iota \in I\) , let \(A_{\iota}\) be a \(\mu_{\iota}\) -measurable subset of \(T_{\iota}\) . For \(A = \prod_{\iota \in I} A_{\iota}\) to be \(\mu\) -measurable, it is necessary and sufficient that one
\end{enumerate}

be in one of the following two cases: \(1^{\circ}\prod_{\iota \in I}\mu_{\iota}(A_{\iota}) = 0;2^{\circ}A_{\iota}\) contains the support of the

measure \(\mu_{\iota}\) , except perhaps for the indices \(\iota\) of a countable subset of I. In each of these two cases, \(\mu(A) = \prod_{\iota \in I} \mu_{\iota}(A_{\iota})\) . (Assuming that one is not in either of the two preceding

cases, show that one can reduce to the case that \(\mu_{\iota}(\mathrm{A}_{\iota}) = 1\) for all \(\iota \in \mathrm{I}\) . Using \(b\) , show that neither A nor T-A can contain a compact set of nonzero measure.)

¶ 14) Let K be the interval {[}0,1{]} of R, λ the Lebesgue measure on K, A an uncountable set, and μ the measure on T = K\^{}A that is the product of the measure λ on each of the factors.

\begin{enumerate}
\def\labelenumi{\alph{enumi})}
\item
  Show that every closed and metrizable subspace X of T is \(\mu\) -negligible (make use of Exer. 13 b) above, and Exer. 8 of GT, IX, §2).
\item
  Let Y be a locally compact space admitting a countable base, \(\nu\) a positive measure on Y, and \(\pi\) a \(\nu\) -proper mapping of Y into T. Show that the image \(\pi(\nu)\) is alien to \(\mu\) (using a), show that \(\pi(\nu)\) is concentrated on a \(\mu\) -negligible set).
\end{enumerate}

\begin{enumerate}
\def\labelenumi{\arabic{enumi})}
\setcounter{enumi}{14}
\tightlist
\item
  Let \((\mathrm{T}_{\iota})_{\iota \in I}\) be any family of compact spaces and, for each \(\iota \in I\) , let \(\mu_{\iota}\) be a positive measure on \(\mathrm{T}_{\iota}\) of total mass 1; let \(\mu\) be the product measure \(\bigotimes_{\iota \in I} \mu_{\iota}\) on \(\mathrm{T} = \prod_{\iota \in \mathrm{I}}\mathrm{T}_{\iota}\) . For every partition (L, M) of I into two subsets, one identifies T with the
\end{enumerate}

product \(\mathrm{T_L} \times \mathrm{T_M}\) , where \(\mathrm{T_L} = \prod_{\iota \in \mathrm{L}} \mathrm{T}_{\iota}\) ; for every \(t \in \mathrm{T}\) , one writes \(t_{\mathrm{L}} = \mathrm{pr}_{\mathrm{L}} t\) , so that

\(t\) is identified with \((t_{\mathrm{L}}, t_{\mathrm{M}})\) ; let \(\mu_{\mathrm{L}}\) be the measure \(\bigotimes_{\iota \in \mathrm{L}} \mu_{\iota}\) on \(\mathrm{T}_{\mathrm{L}}\) . Let \(p\) be a finite

real number \(\geqslant 1\) ; for every function \(\mathbf{f}\) in \(\mathcal{L}_{\mathrm{F}}^{p}(\mathrm{T},\mu)\) (F a Banach space or \(\mathrm{F} = \overline{\mathbf{R}}\) ), one writes \(\mathbf{f}_{\mathrm{L}}(t) = \int \mathbf{f}(t_{\mathrm{L}},t_{\mathrm{M}}) d\mu_{\mathrm{M}}(t_{\mathrm{M}})\) . Show that, with respect to the directed set of finite subsets L of I, \(\mathbf{f}_{\mathrm{L}}\) tends in mean of order \(p\) to \(\mathbf{f}\) , and \(\mathbf{f}_{\mathrm{M}}\) tends in mean of order \(p\) to the constant function equal to \(\int \mathbf{f} d\mu\) . (Approximate \(\mathbf{f}\) by a continuous function depending only on a finite number of variables.)

From this, deduce that if a \(\mu\) -measurable subset \(A\) of \(T\) is such that, for all \(t \in A\) , every point \(t' \in T\) , whose coordinates are equal to those of \(t\) except for a finite number of indices, also belongs to \(A\) , then \(\mu(A) = 0\) or \(\mu(A) = 1\) .

¶ 16) Let \((T_n)\) be an infinite sequence of compact spaces, \(\mu_n\) a positive measure on \(T_n\) of total mass equal to 1, and \(\mu\) the product measure \(\bigotimes_{n=1}^{\infty} \mu_n\) on \(T = \prod_{n=1}^{\infty} T_n\) .

\begin{enumerate}
\def\labelenumi{\alph{enumi})}
\item
  Let \(f\) be a function \(\geqslant 0\) , \(\mu\) -integrable on \(\mathbf{T}\) . Let \((\mathrm{L}_n)\) be an increasing sequence of subsets of \(\mathbf{N}\) , and, setting \(\mathrm{M}_n = \mathbf{N} - \mathrm{L}_n\) , let \(g = \sup_n f_{\mathrm{L}_n}\) , \(h = \sup_n f_{\mathrm{M}_n}\) (the notations of Exer. 15). For every \(\alpha > 0\) , let \(\mathrm{A}_{\alpha}\) be the set of points \(t \in \mathbb{T}\) where \(g(t) > \alpha\) , and \(\mathrm{B}_{\alpha}\) the set of points \(t \in \mathbb{T}\) where \(h(t) > \alpha\) . Show that \(\alpha \cdot \mu(\mathrm{A}_{\alpha}) \leqslant \int f d\mu\) and \(\alpha \cdot \mu(\mathrm{B}_{\alpha}) \leqslant \int f d\mu\) . (Observe that \(\mathrm{A}_{\alpha}\) is the set of \(t \in \mathbb{T}\) where at least one of the \(f_{\mathrm{L}_n}(t)\) is \(> \alpha\) , and express it as a countable union of sets \(\mathrm{G}_n\) , pairwise disjoint and such that \(\alpha \cdot \mu(\mathrm{G}_n) \leqslant \int_{\mathrm{G}_n} f d\mu\) .)
\item
  Suppose that \((\mathrm{L}_n)\) is an increasing sequence of finite subsets of \(\mathbf{N}\) , whose union is \(\mathbf{N}\) . Show that \(f_{\mathrm{L}_n}\) tends almost everywhere to \(f\) , and that \(f_{\mathrm{M}_n}\) tends almost everywhere in \(\mathrm{T}\) to the constant \(\int f d\mu\) . (For every \(\varepsilon > 0\) , consider a continuous function \(g\) depending on only a finite number of variables and such that \(\int |f - g| d\mu \leqslant \varepsilon\) , and apply \(a\) ) to the function \(|f - g|\) .)
\end{enumerate}

\begin{enumerate}
\def\labelenumi{\arabic{enumi})}
\setcounter{enumi}{16}
\tightlist
\item
  Let \((\mathrm{T}_{\iota})_{\iota \in \mathrm{I}}\) be any family of compact spaces; for each \(\iota \in \mathrm{I}\) , let \(\mu_{\iota}\) be a positive measure on \(\mathrm{T}_{\iota}\) , of total mass equal to 1, and let \(\mu\) be the product measure \(\bigotimes_{\iota \in \mathrm{I}} \mu_{\iota}\) on \(\mathrm{T} = \prod_{\iota \in \mathrm{I}} \mathrm{T}_{\iota}\) . For every \(\iota \in \mathrm{I}\) , let \(f_{\iota}\) be a function \(\geqslant 0\) defined on \(\mathrm{T}_{\iota}\) , \(\mu_{\iota}\) -integrable and such that \(\int f_{\iota} d\mu_{\iota} \leqslant 1\) ; set \(\mu_{\iota}' = f_{\iota} \cdot \mu_{\iota}\) and \(\mu' = \bigotimes_{\iota \in \mathrm{I}} \mu_{\iota}'\) (Ch. III, §4, Exer. 6). Assume that \(\mu' \neq 0\) .
\end{enumerate}

\begin{enumerate}
\def\labelenumi{\alph{enumi})}
\tightlist
\item
  For every finite subset \(J\) of \(I\) , set \(T_J = \prod_{\iota \in J} T_\iota\) , \(\mu_J = \bigotimes_{\iota \in J} \mu_\iota\) , \(\mu_J' = \bigotimes_{\iota \in J} \mu_\iota'\) ,
\end{enumerate}

\(f_{\mathrm{J}}(t) = \prod_{\iota \in \mathbb{J}}f_{\iota}(\mathrm{pr}_{\iota}t), g_{\mathrm{J}} = \sqrt{f_{\mathrm{J}}}, \mu_{\mathrm{J}}^{\prime \prime} = \sqrt{\mu_{\mathrm{J}}\mu_{\mathrm{J}}^{\prime}}\) (§5, No.~9) and \(\rho (\mu_{\mathrm{J}},\mu_{\mathrm{J}}^{\prime}) = \mu_{\mathrm{J}}^{\prime \prime}(\mathrm{T}_{\mathrm{J}}) =\)

\(\int g_{\mathrm{J}}d\mu_{\mathrm{J}}\) . Show that, for the functions \(g_{\mathrm{J}}\) to converge in the quadratic mean to a function in \(\mathcal{L}^2 (\mathrm{T},\mu)\) , with respect to the directed ordered set of finite subsets of I, it is necessary and sufficient that the product of the numbers \(\rho (\mu_{\iota},\mu_{\iota}^{\prime})\) be convergent in \(\mathbf{R}_{+}^{*}\) , that is, \(>0\) . (For two finite subsets J,L such that \(\mathrm{J}\subset \mathrm{L}\) , evaluate the norm \(\mathrm{N}_2(g_{\mathrm{J}} - g_{\mathrm{L}})\) by means of the \(\rho (\mu_{\iota},\mu_{\iota}^{\prime})\) .) Deduce from this that the functions \(f_{\mathrm{J}}\) then converge in mean to a function \(f\geqslant 0\) such that \(\mu^{\prime} = f\cdot \mu\) .

\begin{enumerate}
\def\labelenumi{\alph{enumi})}
\setcounter{enumi}{1}
\tightlist
\item
  Show that if the product \(\prod_{\iota \in I} \rho(\mu_{\iota}, \mu_{\iota}')\) is zero, then the measures \(\mu\) and \(\mu'\) are
\end{enumerate}

alien (consider the set \(\mathrm{A_J}\) of points \(t\in \mathrm{T}\) such that \(g_{\mathrm{J}}(t) > 1\) , and evaluate the measures \(\mu (\mathrm{A_J})\) and \(\mu '(T - A_J)\) ).

¶ 18) The notations being the same as in Exer. 17, let \(\nu_{\iota}\) be a positive measure on \(\mathrm{T}_{\iota}\) , of total mass equal to 1, and let \(\nu = \bigotimes \nu_{\iota}\) .

\begin{enumerate}
\def\labelenumi{\alph{enumi})}
\item
  Show that if one of the measures \(\nu_{\iota}\) is alien to the measure \(\mu_{\iota}\) with the same index, then \(\nu\) is alien to \(\mu\) .
\item
  For every \(\iota \in I\) , write \(\nu_{\iota} = \mu_{\iota}' + \mu_{\iota}''\) , where \(\mu_{\iota}'\) has base \(\mu_{\iota}\) , and \(\mu_{\iota}''\) is alien to \(\mu_{\iota}\) (§5, Th. 3); suppose that \(\mu_{\iota}' \neq 0\) for all \(\iota \in I\) . Show that, for \(\nu\) not to be alien to \(\mu\) , it is necessary and sufficient that \(\mu_{\iota}'' = 0\) except for a countable family of indices, and that the product \(\mu' = \bigotimes_{\iota \in I} \mu_{\iota}'\) be a measure \(\neq 0\) with base \(\mu\) ; \(\mu'' = \nu - \mu'\) is then
\end{enumerate}

alien to \(\mu\) . (If \(\mu_{\iota}^{\prime\prime} \neq 0\) for an uncountable infinity of indices, show that there exists a number \(\alpha\) such that \(0 < \alpha < 1\) and such that, for a countable infinity of indices \(\iota_{n}\) , one has \(\mu_{\iota_{n}}^{\prime}(T_{\iota_{n}}) \leqslant \alpha\) ; show then that \(\mu\) is alien to \(\nu\) by using \(a\) ); prove the second part of the proposition by using Exer. 17.)

\begin{enumerate}
\def\labelenumi{\arabic{enumi})}
\setcounter{enumi}{18}
\tightlist
\item
  Let \(X\) be a locally compact space, \(Y\) a locally compact space countable at infinity, \(A\) a universally measurable subset of \(X \times Y\) .
\end{enumerate}

\begin{enumerate}
\def\labelenumi{\alph{enumi})}
\item
  For every \(x \in X\) , show that the section \(\mathrm{A}(x)\) is a universally measurable subset of Y. Moreover, for every positive measure \(\mu\) on Y, the function \(x \mapsto \mu^{*}(\mathrm{A}(x))\) is universally measurable on X (make use of the Lebesgue--Fubini theorem).
\item
  Let \(\mu\) be a positive measure on Y such that, for almost every \(y \in Y\) (for \(\mu\) ), the section \(\overline{A}(y)\) of A is countable. Show that the set N of \(x \in X\) such that \(\mu^{*}(A(x)) > 0\) cannot contain any countable compact set.
\end{enumerate}

\begin{enumerate}
\def\labelenumi{\arabic{enumi})}
\setcounter{enumi}{19}
\tightlist
\item
  \begin{enumerate}
  \def\labelenumii{\alph{enumii})}
  \tightlist
  \item
    Let \(\nu\) be a positive measure on the unit circle \(\mathbf{U}: |z| = 1\) in \(\mathbf{C}\) , and let \(\mu\) be the complex measure \(j \cdot \nu\) , where \(j: z \mapsto z\) is the canonical injection of \(\mathbf{U}\) into \(\mathbf{C}\) , so that \(|\mu| = \nu\) . Let \(\rho\) be the inverse image of \(\nu\) under the canonical local homeomorphism \(\varphi: t \mapsto e^{it}\) of \(\mathbf{R}\) onto \(\mathbf{U}\) (§6, No.~6). Let A be a \(\nu\) -measurable subset of \(\mathbf{U}\) ; if \(\mu(A) = |\mu(A)|e^{i\omega}\) , show that \(|\mu(A)| \leqslant \left|\int_{\omega - \pi/2}^{\omega + \pi/2} e^{i\theta} d\rho(\theta)\right|\) .
  \end{enumerate}
\end{enumerate}

\begin{enumerate}
\def\labelenumi{\alph{enumi})}
\setcounter{enumi}{1}
\tightlist
\item
  Deduce from \(a\) ) that, as A runs over the set of \(\mu\) -measurable subsets of \(\mathbf{U}\) , one has
\end{enumerate}

\[
\begin{array}{l} \sup _ {\mathrm{A}} | \mu (\mathrm{A}) | \geqslant \sup _ {\omega} \left| \int_ {\omega - \pi / 2} ^ {\omega + \pi / 2} \mathcal {R} \big (e ^ {i (\theta - \omega)} \big) d \rho (\theta) \right| \\ \geqslant \frac {1}{2 \pi} \int_ {0} ^ {2 \pi} d \omega \int_ {\omega - \pi / 2} ^ {\omega + \pi / 2} \mathcal {R} \big (e ^ {i (\theta - \omega)} \big) d \rho (\theta) \end{array}
\]

and conclude therefrom that

\[
\sup _ {A} | \mu (A) | \geqslant \frac {1}{\pi} \| \mu \|\tag{*}
\]

(evaluate the last integral with the help of the Lebesgue--Fubini theorem).

\begin{enumerate}
\def\labelenumi{\alph{enumi})}
\setcounter{enumi}{2}
\tightlist
\item
  Show that the inequality (*) is true for every bounded complex measure on a locally compact space X, where A runs over the set of \(\mu\) -measurable subsets of X. (Let \(\mu = h \cdot |\mu|\) , where h takes its values in U. Consider the image measure \(\nu = h(|\mu|)\) on U and show that \(h(\mu) = j \cdot \nu\) ; then apply b) to \(h(\mu)\) .)
\end{enumerate}

\begin{enumerate}
\def\labelenumi{\arabic{enumi})}
\setcounter{enumi}{20}
\tightlist
\item
  Denote by \(\lambda\) the Lebesgue measure on \(\mathbf{R}^n\) .
\end{enumerate}

\begin{enumerate}
\def\labelenumi{\alph{enumi})}
\tightlist
\item
  Let \(\mathbf{C}_0\) be a cube (GT, VI, §1, No.~1) in \(\mathbf{R}^n\) , and let \(u\) be a numerical function on \(\mathbf{C}_0\) that is integrable (for \(\lambda\) ). Let \(s\) be a real number such that
\end{enumerate}

\[
s \geqslant \frac {1}{\lambda (\mathrm{C} _ {0})} \int_ {\mathrm{C} _ {0}} | u | d \lambda .
\]

Show that there exists a sequence \((\mathrm{I}_k)\) of open cubes contained in \(\mathbf{C}_0\) , pairwise disjoint, and satisfying the following conditions:

\(1^{\circ}\left|u(x)\right|\leqslant s\) almost everywhere in the set \(\mathrm{C_0 - }\bigcup_{k}\mathrm{I}_k;\)

\(2^{\circ}\) if one sets

\[
u _ {k} = \frac {1}{\lambda (\mathrm{I} _ {k})} \int_ {\mathrm{I} _ {k}} | u (x) | d \lambda (x),
\]

then \(u_{k} \leqslant 2^{n}s\) for all \(k\) ;

\[
3 ^ {\circ} \sum_ {k} \lambda (I _ {k}) \leqslant \frac {1}{s} \int_ {C _ {0}} | u (x) | d \lambda (x).
\]

(Dividing into two equal intervals each of the projections of \(C_{0}\) , among the \(2^{n}\) pairwise disjoint open cubes choose the cubes \(I_{1i}\) such that the corresponding `mean values' \(u_{1i}\) are all \(\geqslant s\) , and, because of the choice of s, show that \(u_{1i} \leqslant 2^{n}s\) . Similarly decompose into \(2^{n}\) cubes each of the cubes of the first subdivision that does not figure among the \(I_{1i}\) , which yields a second finite family of cubes \(I_{2i}\) chosen among all of the new cubes obtained as those for which the corresponding mean value \(u_{2i}\) is \(\leqslant s\) . Continue by recursion. To prove that the condition \(1^{\circ}\) is verified, argue by contradiction, noting that the subset \(B_{m}\) of \(C_{0} - \bigcup_{k} I_{k}\) where \(|u(x)| \geqslant s + 1/m\) is, up to a negligible set, the intersection of a decreasing sequence \((\mathrm{U}_r)\) of open sets, where \(\mathrm{U}_r\) is the union of the cubes of the \(r\) -th subdivision that do not figure among the \(\mathrm{I}_k\) and that contain a point of \(\mathrm{B}_m\) .

\begin{enumerate}
\def\labelenumi{\alph{enumi})}
\setcounter{enumi}{1}
\tightlist
\item
  Let \(\mathcal{F}\) be the set of integrable numerical functions \(u\) on \(C_0\) having the following property: for every cube \(C \subset C_0\) , if \(m_C(u)\) is the mean value
\end{enumerate}

\[
\frac {1}{\lambda (\mathrm{C})} \int_ {\mathrm{C}} u (x) d \lambda (x),
\]

then \(\int_{\mathbb{C}} |u(x) - m_{\mathbb{C}}(u)| d\lambda(x) \leqslant \lambda(\mathbb{C})\) ; if \(u \in \mathcal{F}\) , the same is true of \(u - c\) for every real number \(c\) . For every function \(u \in \mathcal{F}\) , every number \(\sigma > 0\) and every cube \(\mathbb{C} \subset \mathbb{C}_0\) , denote by \(\mathrm{S}(\sigma, u, \mathbb{C})\) the set of \(x \in \mathbb{C}\) such that \(|u(x) - m_{\mathbb{C}}(u)| \geqslant \sigma\) . Finally, denote by \(\mathrm{G}(\sigma)\) the smallest number such that

\[
\lambda \bigl (\mathrm{S} (\sigma , u, \mathrm{C}) \bigr) \leqslant \mathrm{G} (\sigma) \int_ {\mathrm{C}} | u (x) - m _ {\mathrm{C}} (u) | d \lambda (x)
\]

as \(u\) runs over \(\mathcal{F}\) , and \(C\) over the set of cubes contained in \(C_0\) ; one has \(G(\sigma) \leqslant 1 / \sigma\) . Show that if \(\sigma / 2^n > s \geqslant 1\) then

\[
\mathrm{G} (\sigma) \leqslant \frac {1}{s} \mathrm{G} (\sigma - 2 ^ {n} s).
\]

(One can restrict oneself to proving that, for a \(u \in \mathcal{F}\) , one has

\[
\lambda \left(\mathrm{S} (\sigma , u, \mathrm{C} _ {0})\right) \leqslant \frac {1}{s} \mathrm{G} (\sigma - 2 ^ {n} s) \int_ {\mathrm{C} _ {0}} | u (x) - m _ {\mathrm{C} _ {0}} (u) | d \lambda (x),
\]

and one can suppose, to simplify, that \(m_{\mathrm{C}_0}(u) = 0\) . Make use then of the decomposition of \(\mathrm{C}_0\) into a union of cubes \(\mathrm{I}_k\) and a set where \(|u(x)| \leqslant s\) almost everywhere, on noting that if \(x \in \mathrm{I}_k\) belongs to \(\mathrm{S}(\sigma, u, \mathrm{C}_0)\) , then \(|u(x) - m_{\mathrm{I}_k}(u)| \geqslant \sigma - 2^n s\) and consequently

\[
\lambda \left(\mathrm{I} _ {k} \cap \mathrm{S} (\sigma , u, \mathrm{C} _ {0})\right) \leqslant \mathrm{G} (\sigma - 2 ^ {n} s) \int_ {\mathrm{I} _ {k}} | u (x) - m _ {\mathrm{I} _ {k}} (u) | d \lambda (x). \Big)
\]

\begin{enumerate}
\def\labelenumi{\alph{enumi})}
\setcounter{enumi}{2}
\tightlist
\item
  Deduce from \(b\) ) that there exist two constants \(B, b\) depending only on \(n\) , such that, for every function \(u \in \mathcal{F}\) and every number \(\sigma > 0\) , one has \(\lambda(\mathrm{S}(\sigma, u, \mathrm{C}_0)) \leqslant \mathrm{Be}^{-b\sigma}\lambda(\mathrm{C}_0)\) . (Take \(s = e\) in \(b\) ).)
\end{enumerate}

\chapter*{HISTORICAL NOTE}
\phantomsection
\addcontentsline{toc}{chapter}{HISTORICAL NOTE}

\begin{center}
\textit{(Chapters II to V)}
\end{center}

(N.B. --- The Roman numerals refer to the bibliography at the end of this note.)

The development of the modern concept of integral is closely tied to the evolution of the idea of function, and to the deeper study of numerical functions of real variables, pursued since the beginning of the 19th century. It is known that Euler already conceived the notion of function in a quite general way, since for him the giving of an `arbitrary' curve intersected in a single point by every line parallel to the Oy axis defines a function \(y = f(x)\) (cf.~FRV, Hist. Note for Chs. I--III); however, like most of his contemporaries, he refused to admit that such functions could be expressed `analytically'. This point of view was scarcely to change until the work of Fourier; but the discovery, by the latter, of the possibility of representing discontinuous functions as sums of trigonometric series, \(^{*}\) was to exercise a decisive influence on later generations. To tell the truth, Fourier's proofs were completely lacking in rigor, and their domain of validity was not clearly apparent; however, the integral formulas

\[
a _ {n} = \frac {1}{\pi} \int_ {- \pi} ^ {+ \pi} \varphi (x) \cos n x d x, b _ {n} = \frac {1}{\pi} \int_ {- \pi} ^ {+ \pi} \varphi (x) \sin n x d x (n \geqslant 1),\tag{1}
\]

giving the coefficients of the Fourier series expansion of \(\varphi\) , had an obvious intuitive meaning when \(\varphi\) was assumed to be piecewise continuous and monotone.* Also, it is these functions to which Dirichlet restricts himself at the outset, in the famous memoir (II) in which he established the convergence of the Fourier series; but already, at the end of his work, he is preoccupied with the extension of his results to more extensive classes of functions. One knows that it is on this occasion that Dirichlet, making precise the ideas of Fourier, defined the general concept of function as we understand it today; naturally, the first point to clarify was to know in which cases it was again possible to attach a meaning to the formulas (1). ``When the solutions of continuity {[}of \(\varphi\) {]} are infinite in number \ldots'' says Dirichlet ((II), p.~169), ``it is then necessary that the function \(\varphi(x)\) be such that, if a and b denote any two quantities between \(-\pi\) and \(+\pi\) , one can always place between a and b other quantities r and s sufficiently close that the function remains continuous in the interval from r to s. One easily senses the need for this restriction on considering that the various terms of the {[}Fourier{]} series are definite integrals and on going back to the fundamental concept of integral. It will then be seen that the integral of a function has meaning only insofar as the function satisfies the condition just stated.''

In modern terms, Dirichlet seems to believe that integrability is equivalent to the fact that the points of discontinuity form a nowhere dense set; he notes moreover, a few lines later, the famous example of the function equal to c for x rational, and to a different value d for x irrational, and he asserts that this function ``cannot be substituted'' into the integral. He also announced future works on this subject, but these works were never published,** and for 25 years no one seems to have sought to advance in that direction, perhaps because the consideration of functions so `pathological' appeared at the time to be entirely devoid of interest; at any rate, when Riemann, in 1854 ((III), pp.~227--264), took up the question again (still in connection with trigonometric series**), he felt the need to justify his work: ``Whatever may be our ignorance relating to the manner in which forces and the states of matter vary with time and place in the infinitely small, we can nevertheless take it as certain that the functions to which Dirichlet's researches do not apply, do not arise in natural phenomena. However,'' he continued, ``it seems that these cases not treated by Dirichlet merit attention for two reasons. First, as Dirichlet himself remarks at the end of his work, this subject is very closely related to the principles of the infinitesimal calculus, and may serve to bring greater clarity and certainty to these principles. From this point of view, its study has immediate interest. Secondly, the application of Fourier series is not limited to research in Physics; they are at present also applied with success in a domain of pure mathematics, the theory of numbers, and it appears that it is precisely the functions whose expansion in trigonometric series was not studied by Dirichlet, that show some importance there ((III), pp.~237--238).

The idea of Riemann is to begin from the approximation procedure of the integral, restored to a place of honor by Cauchy, and to determine when the `Riemann sums' of a function \(f\) , in a bounded interval \([a, b]\) , tend to a limit (as the maximal length of the intervals of the subdivision tends to 0), a problem whose solution he obtains without difficulty, in the following form: for every \(\alpha > 0\) , there exists a subdivision of \([a, b]\) into subintervals of maximal length sufficiently small that the sum of the lengths of the intervals of this subdivision on which the oscillation of \(f\) is \(>\alpha\) , is arbitrarily small. He shows, moreover, that this condition is verified, not only for functions piecewise continuous and monotone, but also for some functions that may have a dense set of points of discontinuity.*

Riemann's memoir was not published until after his death, in 1867. This time, however, the epoch was more favorable to this type of research, and the `Riemann integral' took its place naturally in the current of ideas then leading to an intensive study of the `continuum' and of functions of real variables (Weierstrass, Du Bois-Reymond, Hankel, Dini) and which was to conclude, with Cantor, in the birth of the theory of sets. The form given by Riemann to the condition for integrability suggested the idea of a `measure' for the set of points of discontinuity of a function in an interval; but nearly 30 years were to pass before one succeeded in giving a fruitful and convenient definition of this concept.

The first attempts in this direction are due to Stolz, Harnack and Cantor (1884--85); to define the `measure' of a bounded subset E of R, the first two consider the sets \(F \supset E\) that are a finite union of intervals, take for each F the sum of the lengths of the corresponding intervals, and call `measure' of E the infimum of these numbers; whereas Cantor, situating himself right away in the space \(R^{n}\) , considers for a bounded set E and for \(\rho > 0\) the neighborhood \(V(\rho)\) of E formed by the points whose distance to E is \(\leqslant \rho\) , and takes the infimum of the `volume' of \(V(\rho)\) .** With this definition, the `measure' of a set was equal to that of its closure, from which it follows in particular that the `measure' of the union of two disjoint sets may be less than the sum of the `measures' of these two sets. Undoubtedly to alleviate this last difficulty, Peano (V) and Jordan (VI) introduced, several years later, alongside Cantor's `measure' \(\mu(\mathrm{A})\) of a set A contained in a box I, its `inner measure' \(\mu(\mathrm{I}) - \mu(\mathrm{I} - \mathrm{A})\) , and called `measurable' the sets A (nowadays said to be `quarriable') for which these two numbers coincide. The union of two disjoint quarriable sets A, B is then quarriable and has as `measure' the sum of the `measures' of A and B; however, a bounded open set is not necessarily quarriable, nor is the set of rational numbers contained in a bounded interval, which deprived the Peano--Jordan concept of much of its interest.

It is to E. Borel (IX) that is due the merit of discerning the deficiencies of previous definitions and seeing how they can be remedied. It had been known since Cantor that every open set U in R is the union of the countable family of its `components', which are pairwise disjoint open intervals; instead of seeking to approach U `from the outside' by enclosing it in a finite sequence of intervals, Borel, basing himself on the preceding result, proposes to take as the measure of U (when U is bounded) the sum of the lengths of its components. He then describes very summarily* the class of sets (subsequently called `Borel sets') obtainable, starting from the open sets, by indefinitely iterating the operations of countable union and the `difference' A -- B, and he indicates that for these sets, one can define a measure that possesses the fundamental property of complete additivity: if a sequence \((\mathrm{A}_{n})\) consists of pairwise disjoint Borel sets, then the measure of their union (assumed to be bounded) is equal to the sum of their measures.

This definition was to inaugurate a new era in Analysis: on the one hand, combined with the contemporary work of Baire, it formed the point of departure of a whole series of investigations of a topological nature on the classification of sets of points; and, above all, it was to serve as a basis for the extension of the notion of integral, carried out by Lebesgue in the first years of the 20th century.

In his thesis (X a)), Lebesgue begins by making precise and developing the succinct indications of E. Borel; imitating the Peano--Jordan method, the `outer measure' of a bounded set \(A \subset R\) is defined as the infimum of the measures of the open sets containing A; then, if I is an interval containing A, the `inner measure' of A is the difference of the outer measures of I and I - A; one thus obtains a notion of `measurable set', which differs from Borel's original `constructive' definition only by the adjunction of a subset of a set of measure zero in the sense of Borel. This definition extended at once to the spaces \(R^{n}\) ; the old conception of the definite integral \(\int_{a}^{b} f(t) dt\) of a bounded function \(\geqslant 0\) as the `area' bounded by the curve \(y = f(x)\) and the lines x = a, x = b and y = 0, thus provided an immediate extension of the Riemann integral to all of the functions f for which the measure of the preceding set happens to be defined. But the originality of Lebesgue does not reside so much in the idea of this extension* as in his discovery of the fundamental theorem on passage to the limit in the integral so conceived, a theorem that appears in his work as a consequence of the complete additivity of the measure;** he perceives at once its full importance, and makes it the cornerstone of the didactic exposition of his theory that he gives, already in 1904, in his famous ``Leçons sur l'intégration et la recherche des fonctions primitives'' (X b)).***

We cannot describe here in detail the countless advances that Lebesgue's results were to lead to in the study of the classical problems of the infinitesimal calculus; we shall have occasion to stress some of them in later Books. Lebesgue himself, already in his thesis, had applied his theory to extending the classical notions of length and area to sets more general than the usual curves and surfaces; for the considerable development of this theory during the last half-century, we refer the reader to the recent exposition by L. Cesari (XXVI). We mention also the applications to trigonometric series, developed by Lebesgue almost immediately after his thesis (X c), which were to open new horizons to this theory, whose exploration is to this day far from being completed (see (XXV)). Finally, and above all, the definition of the \(L^{p}\) spaces and the Riesz--Fischer theorem ((XIII), (XV a) and (XV b)); cf.~the Historical Note for Book V) brought to light the role that the new concept of integral could play in Functional analysis; a role that was only to increase with the subsequent generalizations of this concept, of which we shall speak in a moment.

But first, let us dwell a little longer on one of the problems to which Lebesgue devoted the most effort, the connection between the concepts of integral and primitive. With the generalization of integral introduced by Riemann, the question naturally arose as to whether the classical correspondence between integral and primitive, valid for continuous functions, again subsisted in more general cases. Now, it is easy to give examples of functions \(f\) , integrable in the sense of Riemann, such that \(\int_{a}^{x} f(t) dt\) has no derivative (not even a right-derivative or a left-derivative) at certain points (cf.~FRV, II, §2, Exer. 1); conversely, Volterra had shown, in 1881, that a function \(F(x)\) can have a derivative that is bounded in an interval I but is not integrable (in the sense of Riemann) in I. By an analysis of great subtlety (in which the theorem on passage to the limit in the integral is far from sufficient), Lebesgue succeeded in showing that if \(f\) is integrable (in his sense) on \([a, b]\) , then \(F(x) = \int_{a}^{x} f(t) dt\) has almost everywhere a derivative equal to \(f(x) (\mathrm{X}b)\) ). Conversely, if a function \(g\) is differentiable on \([a, b]\) and its derivative \(g' = f\) is bounded, then \(f\) is integrable and the formula \(g(x) - g(a) = \int_{a}^{x} f(t) dt\) holds. However, Lebesgue ascertained that the problem is much more complex when \(g'\) is not bounded; in this case \(g'\) is not necessarily integrable, thus the first problem was to characterize the continuous functions \(g\) for which \(g'\) exists almost everywhere and is integrable. Restricting himself to the case that one of the derivatives* of \(g\) is everywhere finite, Lebesgue showed that \(g\) is necessarily a function of bounded variation.** Finally, he established a converse of the last result: a function \(g\) of bounded variation admits a derivative almost everywhere, and \(g'\) is integrable; however, it is no longer necessarily the case that

\[
g (x) - g (a) = \int_ {a} ^ {x} g ^ {\prime} (t) d t;\tag{2}
\]

the difference between the two members of this relation is a function of bounded variation that is nonconstant and has derivative zero almost everywhere (a `singular' function). It remained to characterize the functions g of bounded variation such that relation (2) holds. Lebesgue established that these functions (called `absolutely continuous' by Vitali, who made a detailed study of them) are those having the following property: the total variation of g on an open set U (the sum of the total variations of g on each of the connected components of U) tends to 0 with the measure of U.

We shall see below how these results, in a weakened form, were to acquire later on a much more general significance. In their original form, their field of application has remained rather restricted, not going beyond the framework of the `fine' theory of functions of real variables; they too remain outside the scope of the present treatise.* All the more reason that the same is true of the later developments in the theory of primitives; we shall content ourselves with mentioning here the profound work of Denjoy and his emulators and continuers of his work (Perron, de la Vallée-Poussin, Khintchine, Lusin, Banach, etc.); the reader will find a detailed exposition of these matters in the book of S. Saks (XXIV).

One of the essential advances brought forth by Lebesgue's theory concerns multiple integrals. This concept was introduced towards the middle of the 18th century, initially in `indefinite integral' form (by analogy with the theory of the integral of functions of a single variable, \(\iint f(x,y) dx dy\) denotes a solution of the equation \(\frac{\partial^2 z}{\partial x \partial y} = f(x,y)\) ); but Euler, as early as 1770, had a very clear conception of the double integral extended to a bounded domain (limited by analytic arcs of curves), and he wrote correctly the formula evaluating such an integral by means of two successive single integrals (I). It was not difficult to justify this formula in terms of `Riemann sums', as long as the function being integrated was continuous, and the domain of integration not too complicated; but as soon as one wanted to take up more general cases, the Riemann procedure met serious difficulties ( \(f(x,y)\) ) could be integrable in the sense of Riemann, without \(\int dx \int f(x,y) dy\) having meaning when the single integrals are taken in the sense of Riemann). These difficulties vanished when one passed to the definition of Lebesgue; the latter had already shown in his thesis that when \(f(x,y)\) is a bounded `Baire function', so are the functions \(y \mapsto f(x,y)\) (for every \(x\) ) and \(x \mapsto \int f(x,y) dy\) , and one has the formula

\[
\iint f (x, y) d x d y = \int d x \int f (x, y) d y\tag{3}
\]

(the integral being taken over a rectangle). A little later, Fubini (XIV) brought to this result an important complement by proving that if one only assumes f to be integrable, then the set of x such that \(y \mapsto f(x, y)\) is not integrable is of measure zero, which permitted immediately extending the formula (3) to this case.

Finally, in 1910 (X d)), Lebesgue undertakes the extension, to multiple integrals, of his results on the derivatives of single integrals. He is thus led to associate to a function \(f\) , integrable on every compact subset of \(\mathbf{R}^n\) , the set function \(\mathrm{F(E)} = \int_{\mathrm{E}} f(\mathbf{x}) \, dx\) , defined for every integrable subset \(\mathrm{E}\) of \(\mathbf{R}^n\) , which generalizes the concept of `indefinite integral'; and he observes on this occasion that this function has the following two properties: \(1^\circ\) it is completely additive; \(2^\circ\) it is `absolutely continuous' in the sense that \(\mathrm{F(E)}\) tends to 0 with the measure of \(\mathrm{E}\) . The essential part of Lebesgue's memoir consists in proving the converse of this proposition.* But he does not rest there, and, in the same direction, calls attention to the possibility of generalizing the concept of function of bounded variation, by considering functions \(\mathrm{F(E)}\) of a measurable set, completely additive and such that \(\sum_{n} |\mathrm{F(E}_{n})|\) remains bounded for every countable partition of \(\mathrm{E}\) into measurable sets \(\mathrm{E}_{n}\) . And, if he in fact restricts himself to considering such functions only on the set of boxes of \(\mathbf{R}^n\) , it is quite clear that there remained only one more step to arrive at the general notion of measure that J. Radon was to define in 1913, encompassing in a single synthesis the Lebesgue integral and the Stieltjes integral, of which we must now speak.

In 1894, T. Stieltjes published, under the title ``Recherches sur les fractions continues'' (VIII), a highly original memoir in which, starting from a seemingly quite special question, problems of an entirely novel nature in the theory of analytic functions and functions of a real variable were posed and solved, with rare elegance. \(^{**}\) In order to represent the limit of a certain sequence of analytic functions, Stieltjes was led, among other things, to introduce the concept of a positive `distribution of mass' on the line, a concept long familiar in the physical sciences but which had only been considered theretofore in mathematics under restrictive hypotheses (in general, the existence of a `density' at every point, varying in a continuous manner); he observes that the giving of such a distribution is equivalent to that of the increasing function \(\varphi(x)\) that gives the total mass contained in the interval with end-points 0 and x for x \textgreater{} 0, and this mass with the sign changed for x < 0, the discontinuities of \(\varphi\) corresponding to the masses `concentrated at a point'.* Stieltjes then forms, for such a mass distribution in an interval \([a,b]\) , the `Riemann sums' \(\sum_{i} f(\xi_i)(\varphi(x_{i+1}) - \varphi(x_i))\) and shows that, when \(f\) is continuous on \([a,b]\) , these sums tend to a limit that he denotes \(\int_{a}^{b} f(x) d\varphi(x)\) . Having no need to integrate functions other than continuous (or even differentiable) ones, Stieltjes did not push further the study of this integral** and for a decade the concept seems not to have attracted any attention.*** However, in 1909, F. Riesz (XV c)), solving a problem posed several years earlier by Hadamard (cf.~the Hist. Note of Book V), proved that the Stieltjes integrals \(f \mapsto \int_{a}^{b} f d\varphi\) are the most general continuous linear functionals on the space \(\mathcal{C}(\mathrm{I})\) of continuous real-valued functions on \(\mathrm{I} = [a,b]\) ( \(\mathcal{C}(\mathrm{I})\) being equipped with the topology of uniform convergence);**** and the elegance and simplicity of this result almost immediately brought forth various generalizations. The most felicitous was that of J. Radon, in 1913 (XVII): combining the ideas of F. Riesz and Lebesgue, he showed how one could define an integral by Lebesgue's procedures, starting from an arbitrary `completely additive set function' (defined on the measurable sets for Lebesgue measure) instead of starting from Lebesgue measure. In the concept of `Radon measure' on \(\mathbf{R}^n\) thus defined, one found absorbed that of a function `of bounded variation': the decomposition of such a function as the difference of two increasing functions is a special case of the decomposition of a measure as the difference of two positive measures; similarly, a `measure with base \(\mu'\) corresponds to the notion of 'absolutely continuous' function, and the decomposition of an arbitrary measure into a measure with base \(\mu\) and a measure alien to \(\mu\) , to the Lebesgue decomposition of a function of bounded variation as a sum of an absolutely continuous function and a `singular' function. Moreover, Radon showed that the `density' with respect to \(\mu\) of a measure with base \(\mu\) again exists when \(\mu\) is a measure having Lebesgue measure as base, by using an earlier idea of F. Riesz (taken up again later and popularized by J. von Neumann, among others), which consists in constructing an image of the measure \(\mu\) under a mapping \(\theta\) of \(\mathbf{R}^n\) into \(\mathbf{R}\) , chosen so that \(\theta(\mu)\) is Lebesgue measure on \(\mathbf{R}\) (cf.~Ch. V, §6, Exer. 8 c)).

Almost immediately after the appearance of Radon's memoir, Fréchet observed that nearly all of the results in that work could be extended to the case where the `completely additive set function', instead of being defined for the measurable subsets of \(\mathbf{R}^n\) , is defined for certain subsets of an arbitrary set E (these subsets being such that the operations of countable union and of `difference' yield sets for which the function is again defined). However, the expression of a measure with base \(\mu\) in the form \(g \cdot \mu\) rested, with Lebesgue and Radon, on arguments involving the topology of \(\mathbf{R}^n\) in an essential way (and we have seen that Radon's proof is only applicable if \(\mu\) is a measure having Lebesgue measure as base); it was not until 1930 that O. Nikodym (XX) obtained the theorem in its general form, by a direct argument (notably simplified several years later by J. von Neumann, thanks to the use of properties of \(\mathrm{L}^2\) spaces ((XXII), pp.~127-130)).

With Radon's memoir, the general theory of integration may be considered as completed in its broad outlines; as later substantial acquisitions, one can only scarcely mention the definition of an infinite product of measures, due to Daniell (XIX \(b\) ), and that of the integral of a function with values in a Banach space, given by Bochner in 1933 (XXI), which paved the way for the study of the `weak integral' which we shall treat in Chapter VI. But it remained to popularize the new theory, and to make of it a mathematical instrument for everyday use, when the majority of mathematicians, around 1910, as yet viewed the `Lebesgue integral' only as an instrument of great precision, delicate to manipulate, destined solely for research of extreme subtlety and extreme abstraction. That was to be the work of Carathéodory, in a book that has long remained a classic (XVIII) and which, moreover, enriched the theory of Radon with numerous original observations.

But it is also with this book that the concept of integral, which had been in the forefront of Lebesgue's preoccupations (amply indicated by the titles of his thesis (X a)) and his principal work on these questions (X b))), for the first time gave way to that of measure, which had been an auxiliary technique with Lebesgue (as with Jordan before him). This change in point of view was undoubtedly due, for Carathéodory, to the excessive importance that he seemed to have attached to `p-dimensional measures'.* Ever since, authors who have treated integration have been divided between these two points of view, not without entering into debates that have caused the flow of much ink, if not much blood.* Some have followed Carathéodory; in their ever more abstract and ever more axiomatized expositions, the measure, with all the technical refinements to which it lends itself, not only plays the dominant role but also tends to lose contact with the topological structures to which it is in fact tied in most of the problems in which it arises. Other expositions, such as the present treatise, follow more or less closely a method already indicated in 1911 by W. H. Young, in a memoir unfortunately little noticed (XI b)), and subsequently developed by Daniell. The former, dealing with the Lebesgue integral, started from the `Cauchy integral' of continuous functions with compact support, assumed to be known, in order to define successively (as we have done in Ch. IV, §1) the upper integral of lower semi-continuous functions, then of arbitrary numerical functions, whence a definition of the integrable functions, patterned after that of Lebesgue for sets, by purely `functional' means. Daniell, in 1918 ((XIX a)), cf.~(XXVII)) extended this exposition, with several variants, to functions defined on an arbitrary set; his principal merit was in perceiving the role played in the abstract theory by the condition \(\lim_{n\to \infty}\int f_n d\mu = 0\) for every decreasing sequence \((f_n)\) tending pointwise to 0 (which could not appear as clearly in the theory of Radon measures, where this condition is automatically satisfied by virtue of Dini's theorem). In the same order of ideas (and in strict relation to the methods used in spectral theory prior to Gelfand), we must also call attention to the memoir of F. Riesz (XV d)) that puts into a concise and elegant form the several results from the theory of ordered spaces that play a role in the theory of integration; we have followed his exposition quite closely in Ch. II.

Rather than in expository works, more or less pleasant to read, but whose essential content can no longer vary much, it is on the side of applications that one must search for the progress realized by the theory of integration since 1920: the theory of probability (formerly a pretext for puzzles and paradoxes, and having become a branch of the theory of integration since its axiomatisation by Kolmogoroff (XXIII), albeit an autonomous branch with its own methods and problems); ergodic theory; spectral theory and harmonic analysis, since the discovery by Haar of the measure that bears his name, and the movement of ideas provoked by this discovery, have made of the integral one of the most important tools in the theory of groups. With these questions we leave the framework of the present Note; some of them will be treated in later chapters or Books.

\chapter*{BIBLIOGRAPHY}
\phantomsection
\addcontentsline{toc}{chapter}{BIBLIOGRAPHY}

\begin{enumerate}
\def\labelenumi{(\Roman{enumi})}
\item
  L. EULER, Opera Omnia: De formulas integralibus duplicatis (1), v. XVII, Leipzig--Berlin (Teubner), 1915, pp.~289--315.
\item
  G. LEJEUNE--DIRICHLET, Sur la convergence des séries trigonométriques qui servent à représenter une fonction arbitraire entre des limites données, J. de Crelle, 4 (1829), pp.~157--169 (= Werke, v. I, pp.~118--132, Berlin (G. Reimer), 1889).
\item
  B. RIEMANN, Gesammelte Mathematische Werke, 2nd. edn., Leipzig (Teubner), 1892.
\item
  G. CANTOR, Gesammelte Abhandlungen, Berlin (Springer), 1932.
\item
  G. PEANO, Applicazioni geometriche del calcolo infinitesimale, Turin, 1887.
\item
  C. JORDAN, Cours d'Analyse, v. I, 2nd. edn., Paris (Gauthier-Villars), 1893.
\item
  C. ARZELÀ: a) Sulla integrabilità di una serie di funzioni, Rendic. Acc. dei Lincei, (4), 1 (1885), pp.~321--326; b) Sulla integrazione per serie, ibid., pp.~532--537 and 566--569.
\item
  T. STIELTJES, Recherches sur les fractions continues, Ann. Fac. Sci. de Toulouse, 8 (1894), J. 1 to J. 122.
\item
  E. BOREL, Leçons sur la théorie des fonctions, Paris (Gauthier-Villars), 1898.
\item
  H. LEBESGUE: a) Intégrale, longueur, aire, Annali di Mat., (3), 7 (1902), pp.~231--359; b) Leçons sur l'Intégration et la recherche des fonctions primitives, Paris (Gauthiers--Villars), 1904; c) Sur les séries trigonométriques, Ann. Éc. Norm. Sup., (3), 20 (1903), pp.~453--485; d) Sur l'intégration des fonctions discontinues, Ann. Éc. Norm. Sup., (3), 27 (1910), pp.~361--450.
\item
  W. H. YOUNG: a) On upper and lower integration, Proc. Lond. Math. Soc., (2), 2 (1905), pp.~52--66; b) A new method in the theory of integration, Proc. Lond. Math. Soc., (2), 9 (1911), pp.~15--50.
\item
  G. VITALI: a) Una proprieta delle funzioni misurabili, R. Ist. Lombardo, Rendiconti, (2), 38 (1905), pp.~599--603; b) Sui gruppi di punti e sulle funzioni di variabili reali, Rendic. Acc. Sci. di Torino, 43 (1908), pp.~229--236.
\item
  E. FISCHER, Sur la convergence en moyenne, C. R. Acad. Sci., 144 (1907), pp.~1022--1024.
\item
  G. FUBINI, Sugli integrali multipli, Rendic. Acc. dei Lincei, (5), 16 (1907), pp.~608--614.
\item
  F. RIESZ: a) Sur les systèmes orthogonaux de fonctions, C. R. Acad. Sci., 144 (1907), pp.~615--619; b) Untersuchungen über Systeme integrierbarer Funktionen, Math. Ann., 69 (1910), pp.~449--497; c) Sur certains systèmes singuliers d'équations intégrales, Ann. Éc. Norm. Sup., (3), 28 (1911), pp.~33--62; d) Sur quelques notions fondamentales dans la théorie générale des opérations linéaires, Ann. of Math., (2), 41 (1940), pp.~174--206.
\item
  D. EGOROFF, Sur les suites de fonctions mesurables, C. R. Acad. Sci., 152 (1911), p.~244.
\item
  J. RADON, Theorie und Anwendungen der absolut additiven Mengenfunktionen, Sitzungsber. der math. naturwiss. Klasse der Akad. der Wiss. (Wien), 122, Abt. IIa (1913), pp.~1295--1438.
\item
  C. CARATHÉODORY, Vorlesungen über reelle Funktionen, Leipzig--Berlin (Teubner), 1918.
\item
  P. J. DANIELL: a) A general form of integral, Ann. of Math., (2), 19 (1918), pp.~279--294; b) Integrals in an infinite number of dimensions, Ann. of Math., (2), 20 (1919), pp.~281--288.
\item
  O. NIKODYM, Sur une généralisation des intégrales de M. J. Radon, Fund. Math., 15 (1930), pp.~131--179.
\item
  S. BOCHNER, Integration von Funktionen deren Werte die Elemente eines Vektorraumes sind, Fund. Math., 20 (1933), pp.~262--276.
\item
  J. von NEUMANN, On rings of operators. III, Ann. of Math., (2), 41 (1940), pp.~94--161.
\item
  A. KOLMOGOROFF, Grundbegriffe der Wahrscheinlichkeitsrechnung, Berlin (Springer), 1933.
\item
  S. SAKS, Theory of the integral, 2nd. edn., New York (Stechert), 1937.
\item
  A. ZYGMUND, Trigonometric series, Warsaw, 1935 (2nd. edn., Cambridge University Press, 1959).
\item
  L. CESARI, Surface area, Princeton, 1954.
\item
  L. H. LOOMIS, An introduction to abstract harmonic analysis, London--New York--Toronto (van Nostrand), 1953.
\end{enumerate}

\chapter{Vectorial integration}

In this chapter, if F denotes a Hausdorff locally convex vector space (over R or C), we denote by \(F'\) its dual, by \(F''\) its bidual, and by \(F'^*\) the algebraic dual of \(F'\) (the space of all linear forms on \(F'\) ); \(F''\) is a linear subspace of \(F'^*\) , and F may be identified (as a vector space without topology) with a linear subspace of \(F''\) . We denote by \(F_\sigma\) the vector space F equipped with the weakened topology \(\sigma(F, F')\) ; the qualifiers `weak' and `weakly' refer to this topology.

In this chapter, T denotes a locally compact space, \(\mathcal{K}_{\mathbf{R}}(\mathrm{T})\) or \(\mathcal{K}(\mathrm{T})\) (resp. \(\mathcal{K}_{\mathbf{C}}(\mathrm{T})\) ) the vector space of real (resp. complex) functions on T, continuous and with compact support; for every subset A of T, \(\mathcal{K}(\mathrm{T},\mathrm{A})\) (resp. \(\mathcal{K}_{\mathbf{C}}(\mathrm{T},\mathrm{A})\) ) denotes the subspace of \(\mathcal{K}(\mathrm{T})\) (resp. \(\mathcal{K}_{\mathbf{C}}(\mathrm{T})\) ) formed by the functions whose support is contained in A. Absent express mention to the contrary, the space \(\mathcal{K}(\mathrm{T})\) (resp. \(\mathcal{K}_{\mathbf{C}}(\mathrm{T})\) ) will be equipped with the direct limit topology of the topologies of uniform convergence on each of the subspaces \(\mathcal{K}(\mathrm{T},\mathrm{K})\) (resp. \(\mathcal{K}_{\mathbf{C}}(\mathrm{T},\mathrm{K})\) ), where K runs over the set of compact subsets of T.

We recall that this topology is finer than the topology of uniform convergence, hence is Hausdorff; it induces on each \(\mathcal{K}(\mathrm{T},\mathrm{K})\) (resp. \(\mathcal{K}_{\mathbf{C}}(\mathrm{T},\mathrm{K})\) ) the topology of uniform convergence (TVS, II, §4, No.~4, Remark). The space \(\mathcal{K}_{\mathbf{C}}(\mathrm{T})\) may be identified with the space obtained from \(\mathcal{K}(\mathrm{T})\) by extension of the scalars from R to C. To say that a linear form on \(\mathcal{K}(\mathrm{T})\) is a measure means, by definition, that it is continuous (Ch. III, §1, No.~3, Def. 2).

\section{INTEGRATION OF VECTOR-VALUED FUNCTIONS}

Throughout this section, \(\mu\) denotes a positive measure on T, and F a Hausdorff locally convex vector space over R. For every mapping f of T into F, and every element \(z'\) of the dual \(F'\) of F, we denote by \(\langle f, z' \rangle\) or \(\langle z', f \rangle\) the numerical function \(z' \circ f\) on T. We shall say that f has a property P scalarly if, for every \(z' \in F'\) , \(\langle z', f \rangle\) has the property P. For example, we shall say that f is scalarly essentially \(\mu\) -integrable if, for every \(z' \in F'\) , \(\langle z', f \rangle\) is essentially \(\mu\) -integrable. \(^{1}\)

Note that in this definition, the topology of F intervenes only through the intermediary of the dual \(F'\) of F. If a function f has the property P scalarly, then it again has the property P scalarly when the topology of F is replaced by any Hausdorff locally convex topology compatible with the duality between F and \(F'\) .

\subsection{Scalarly essentially integrable functions}

If \(\mathbf{f}\) is a scalarly essentially \(\mu\) -integrable mapping of \(\mathrm{T}\) into \(\mathrm{F}\) , the mapping \(\mathbf{z}' \mapsto \int \langle \mathbf{f}(t), \mathbf{z}' \rangle d\mu(t)\) is a linear form on \(\mathrm{F}'\) , that is, an element of the algebraic dual \(\mathrm{F}'^*\) .

DEFINITION 1. --- One calls integral of f with respect to \(\mu\) , and denotes by \(\int f d\mu\) or \(\int f(t) d\mu(t)\) , the element of \(F'^{*}\) defined by

\[
\left\langle \mathbf {z} ^ {\prime}, \int \mathbf {f} d \mu \right\rangle = \int \left\langle \mathbf {z} ^ {\prime}, \mathbf {f} \right\rangle d \mu
\]

for all \(\mathbf{z}' \in \mathrm{F}'\) .

If f is continuous with compact support, it is scalarly integrable and Def. 1 coincides with the definition of the integral of f given in Ch. III, §3, No.~1. On the other hand, if F is a Banach space and f is essentially integrable (Ch. V, §1, No.~3, Def. 3), then f is scalarly essentially integrable and Def. 1 coincides with the definition of the integral of f given in Chapter V (Ch. V, §1, No.~3 and Ch. IV, §4, No.~2, Cor. 1 of Th. 1).

Example. --- Let X be a locally compact space, \(t \mapsto \lambda_{t}\) a mapping of T into the space \(\mathcal{M}(X)\) of measures on X. To say that the family \(t \mapsto \lambda_{t}\) is \(\mu\) -adequate means that it consists of positive measures and that the mapping \(t \mapsto \lambda_{t}\) is scalarly essentially \(\mu\) -integrable and \(\mu\) -measurable for the topology \(\sigma(\mathcal{M}(X), \mathcal{K}(X))\) .² Its integral with respect to \(\mu\) is the measure that was denoted \(\int \lambda_t d\mu(t)\) in Ch. V, §3, No.~1.

Remarks. --- 1) If F is finite-dimensional, then every scalarly essentially integrable mapping of T into F is essentially integrable (Ch. V, §1, No.~3). However, in the general case, a scalarly negligible function on a compact space T may even fail to be \(\mu\) -measurable (Exer. 12).

\begin{enumerate}
\def\labelenumi{\arabic{enumi})}
\setcounter{enumi}{1}
\tightlist
\item
  It is clear that the integral of \(\mathbf{f}\) depends only on the class of \(\mathbf{f}\) modulo the space of mappings of \(\mathbf{T}\) into \(\mathbf{F}\) that are scalarly locally \(\mu\) -negligible. Note that a scalarly locally negligible function \(\mathbf{g}\) is not necessarily zero almost everywhere (Exer. 12). However, this is indeed the case when there exists in \(\mathbf{F}'\) a sequence \((\mathbf{z}_n')\) that is dense for the topology \(\sigma(\mathbf{F}',\mathbf{F})\) : for, if \(\mathbf{H}_n\) is the locally negligible set of points \(t \in \mathbf{T}\) such that \(\langle \mathbf{g}(t), \mathbf{z}_n' \rangle \neq 0\) , the union \(\mathbf{H}\) of the \(\mathbf{H}_n\) is locally negligible and, for every \(t \notin \mathbf{H}\) , one has \(\langle \mathbf{g}(t), \mathbf{z}_n' \rangle = 0\) for all \(n\) , whence \(\mathbf{g}(t) = 0\) .
\end{enumerate}

Let \(u\) be a continuous linear mapping of \(F\) into a Hausdorff locally convex space \(G\) ; its transpose \(^{t}u\) is a linear mapping of \(G'\) into \(F'\) , and the (algebraic) transpose \(^{t}(^{t}u)\) is a linear mapping of \(F'^{*}\) into \(G'^{*}\) that extends \(u\) , which we shall again denote by \(u\) . With this convention:

PROPOSITION 1. --- If f is a mapping of T into F that is scalarly essentially μ-integrable, then the mapping u ∘ f is scalarly essentially μ-integrable and

\[
\int (u \circ \mathbf {f}) d \mu = u \Big (\int \mathbf {f} d \mu \Big).
\]

For every \(z^{\prime} \in G^{\prime}\) , \(\langle z^{\prime}, u \circ f \rangle = \langle {}^{t} u(z^{\prime}), f \rangle\) , whence the first assertion; the second follows from the formula

\[
\left\langle \mathbf {z} ^ {\prime}, \int (u \circ \mathbf {f}) d \mu \right\rangle = \int \left\langle \mathbf {z} ^ {\prime}, u \circ \mathbf {f} \right\rangle d \mu = \left\langle {} ^ {t} u (\mathbf {z} ^ {\prime}), \int \mathbf {f} d \mu \right\rangle = \left\langle \mathbf {z} ^ {\prime}, u \left(\int \mathbf {f} d \mu\right) \right\rangle .
\]

In particular, if f is scalarly essentially \(\mu\) -integrable, then it remains scalarly essentially \(\mu\) -integrable when the topology of F is replaced by a coarser topology.

PROPOSITION 2. --- Let f be a scalarly essentially \(\mu\) -integrable mapping of T into F. For every numerical function \(g \geqslant 0\) that is \(\mu\) -measurable and bounded, the mapping \(t \mapsto g(t)\mathbf{f}(t)\) (denoted gf or fg) of T into F is scalarly essentially \(\mu\) -integrable, f is scalarly essentially \((g \cdot \mu)\) -integrable, and

\[
\int \mathbf {f} d (g \cdot \mu) = \int \mathbf {f} g d \mu .
\]

This is an immediate consequence of the formula \(\langle z', gf\rangle = g\langle z', f\rangle\) , for all \(z' \in F'\) , and the formula \(\int h d(g \cdot \mu) = \int hg d\mu\) for every essentially \(\mu\) -integrable scalar function h (Ch. V, §5, No.~3, Th. 1).

A great many propositions about essentially integrable numerical functions may be transposed word for word into propositions about scalarly essentially integrable vector-valued functions. Among the more important, we call attention to the conditions for a function to be essentially integrable with respect to a measure defined by a density (Ch. V, §5, No.~3, Th. 1), or with respect to the image of a measure (Ch. V, §6, No.~2, Th. 1), or with respect to an induced measure (Ch. V, §7, No.~1, Th. 1), or with respect to the sum of a summable family of positive measures (Ch. V, §2, No.~2, Props. 1 and 3 and Cor. 3 of Prop. 1). These transpositions are left to the reader.

However, to obtain statements corresponding to the theorems on `double integrals' (Ch. V, §3, No.~3, Th. 1 and §8, No.~4, Th. 1 (Lebesgue--Fubini theorem)), it is necessary to strengthen the hypotheses (cf.~Exer. 1); on applying the previously cited theorems to each of the functions \(\langle z', f \rangle\) , where \(z' \in F'\) , one thus obtains the following propositions:

PROPOSITION 3. --- Let X be a locally compact space, \(t \mapsto \lambda_{t}\) a \(\mu\) -adequate \(^{3}\) family (Ch. V, §3, No.~1, Def. 1) of positive measures on X, and let \(\nu = \int \lambda_{t} d\mu(t)\) . Let f be a mapping of X into F; assume that \(1^{\circ} f\) is scalarly \(\nu\) -integrable; \(2^{\circ}\) there exists a locally \(\mu\) -negligible set \(N \subset T\) such that, for every \(t \notin N\) , f is scalarly \(\lambda_{t}\) -integrable and \(\int f d\lambda_{t} \in F\) . Then, the function \(t \mapsto \int f d\lambda_{t}\) , defined for \(t \notin N\) , is scalarly essentially \(\mu\) -integrable, and \(^{4}\)

\[
\int \mathbf {f} (x) d \nu (x) = \int d \mu (t) \int \mathbf {f} (x) d \lambda_ {t} (x).
\]

PROPOSITION 4. --- Let T and T' be two locally compact spaces, \(\mu\) (resp. \(\mu'\) ) a positive measure on T (resp. T'), \(\nu = \mu \otimes \mu'\) the product measure on X = T × T'. Let f be a mapping of X into F. Assume that: \(1^{\circ} f\) is scalarly \(\nu\) -integrable; \(2^{\circ}\) there exists a locally \(\mu\) -negligible set N ⊂ T such that, for every \(t \notin N\) , the mapping \(t' \mapsto f(t, t')\) is scalarly \(\mu'\) -integrable, and \(\int f(t, t') d\mu'(t') \in F\) . Then, the function \(t \mapsto \int f(t, t') d\mu'(t')\) , defined for \(t \notin N\) , is scalarly essentially \(\mu\) -integrable and \(^{4}\)

\[
\iint \mathbf {f} (t, t ^ {\prime}) d \mu (t) d \mu^ {\prime} (t ^ {\prime}) = \int d \mu (t) \int \mathbf {f} (t, t ^ {\prime}) d \mu^ {\prime} (t ^ {\prime}).
\]

\subsection{Properties of the integral of a scalarly essentially integrable function}

PROPOSITION 5. --- Let \(\mu\) be a bounded positive measure on T, S a \(\mu\) -measurable set carrying \(\mu\) (Ch. V, §5, No.~7, Def. 4), f a scalarly \(\mu\) -integrable (*) function with values in F. Let D be the closed convex envelope of \(\mathbf{f}(\mathrm{S})\) in the space \(\mathrm{F}'^*\) equipped with the topology \(\sigma(\mathrm{F}'^*, \mathrm{F}')\) . Then \(\int \mathbf{f} d\mu \in \mu(\mathrm{T})\mathrm{D}\) .

Since D is the intersection of the closed half-spaces containing \(\mathbf{f}(\mathrm{S})\) (TVS, II, §5, No.~3, Cor. 1 of Prop. 4), it suffices to prove that the relation \(\langle\mathbf{f}(t),\mathbf{z}^{\prime}\rangle\leqslant a\) for all \(t\in S\) (where \(z^{\prime}\in F^{\prime},a\in R\) ) implies \(\langle z^{\prime},\int f d\mu\rangle\leqslant a\cdot\mu(T)\) ; but since \(\int f d\mu=\int_{S}f d\mu\) , this follows from Ch. IV, §4, No.~2, Cor. 1 of Th. 1.

COROLLARY. --- Let \(\mu\) be a bounded positive measure on T, S a \(\mu\) -measurable set carrying \(\mu\) , and f a mapping of T into F, scalarly \(\mu\) -measurable and such that \(\mathbf{f}(\mathbf{S})\) is contained in a weakly compact convex subset A of F. Then f is scalarly \(\mu\) -integrable, and \(\int f d\mu \in \mu(T)A \subset F\) .

For every \(z^{\prime} \in F^{\prime}\) , \(\langle z^{\prime}, f \rangle\) is \(\mu\) -measurable and bounded in S, hence integrable, which proves that f is scalarly integrable. Moreover, since A is compact in \(F_{\sigma}\) , it is closed in \(F^{\prime*}\) , and the closed convex envelope of \(\mathbf{f}(S)\) in \(F^{\prime*}\) is contained in A, whence the corollary.

PROPOSITION 6. --- Let f be a scalarly essentially μ-integrable function with values in F, such that \(\int f d\mu \in F\) . For every lower semi-continuous semi-norm q on F, 5

\[
q \left(\int \mathbf {f} d \mu\right) \leqslant \int^ {\bullet} (q \circ \mathbf {f}) d \mu .
\]

Let D be the set of \(z \in F\) such that \(q(\mathbf{z}) \leqslant 1\) ; D is closed, convex, and contains 0, therefore \(D = D^{\circ\circ}\) (TVS, II, §6, No.~3, Cor. 3 of Th. 1). It therefore suffices to prove that for every \(z' \in D^{\circ}\) , one has \(\left|\left\langle z', \int f d\mu \right\rangle\right| \leqslant \int^{\bullet}(q \circ f) d\mu\) ; but this follows at once from the fact that \(|\langle z', f(t) \rangle| \leqslant q(f(t))\) for every \(t \in T\) .

Note that the numerical function \(q \circ f\) need not be \(\mu\) -measurable (Exer. 12).

PROPOSITION 7. --- Let f be a mapping of T into F, scalarly essentially \(\mu\) -integrable, such that for every compact subset K of T, \(\mathbf{f}(\mathbf{K})\) is contained in a weakly compact, balanced, convex subset of F. Then \(\int f d\mu\) belongs to the bidual \(F''\) of F.

For every compact subset \(\mathrm{K}\) of \(\mathrm{T}\) ,

\[
\int \mathbf {f} \varphi_ {\mathrm{K}} d \mu = \int (\mathbf {f} \varphi_ {\mathrm{K}}) d (\varphi_ {\mathrm{K}} \cdot \mu);
\]

the Cor. of Prop. 5 can be applied to the bounded measure \(\varphi_{\mathrm{K}} \cdot \mu\) and the function \(\mathbf{f}\varphi_{\mathrm{K}}\) , consequently \(\int \mathbf{f}\varphi_{\mathrm{K}} d\mu \in \mathrm{F}\) . For every \(\mathbf{z}' \in \mathrm{F}'\) , \(\langle \mathbf{z}', \mathbf{f} \rangle\) is essentially \(\mu\) -integrable, consequently (Ch. V, §1, No.~3, Prop. 10)

\[
\int \langle {\bf z} ^ {\prime}, {\bf f} \rangle d \mu = \lim _ {\mathrm{K}} \int \langle {\bf z} ^ {\prime}, {\bf f} \rangle \varphi_ {\mathrm{K}} d \mu ,
\]

the limit being taken with respect to the increasing directed set of compact subsets of T. One concludes that, with respect to this set, \(\int f\varphi_{K}d\mu\) converges to \(\int f d\mu\) for the topology \(\sigma(F'^{*},F')\) . Now,

\[
\left| \left\langle \mathbf {z} ^ {\prime}, \int \mathbf {f} \varphi_ {\mathrm{K}} d \mu \right\rangle \right| = \left| \int \langle \mathbf {z} ^ {\prime}, \mathbf {f} \rangle \varphi_ {\mathrm{K}} d \mu \right| \leqslant \int | \langle \mathbf {z} ^ {\prime}, \mathbf {f} \rangle | d \mu ,
\]

which proves that the set of elements \(\int f\varphi_{K}d\mu\) is a bounded subset of \(F_{\sigma}\) , hence also of F (TVS, IV, §1, No.~1, Prop. 1). Proposition 7 is therefore a consequence of the following lemma:

Lemma 1. --- The closure in \(\mathbf{F}^{\prime*}\) (for the topology \(\sigma(\mathbf{F}^{\prime*},\mathbf{F}^{\prime})\) ) of every bounded subset of \(\mathbf{F}\) is contained in the bidual \(\mathbf{F}''\) .

For, a bounded subset of F is contained in the polar (in \(F''\) ) of a neighborhood of 0 in the strong dual \(F'\) of F, hence is relatively compact in \(F''\) for \(\sigma(\mathrm{F}'',\mathrm{F}')\) (TVS, III, §3, No.~5, Prop. 7 and No.~4, Cor. 2 of Prop. 4); since \(\sigma(\mathrm{F}'',\mathrm{F}')\) is induced by \(\sigma(\mathrm{F}'*,\mathrm{F}')\) , the lemma is proved.

COROLLARY. --- Suppose F is semi-reflexive, and let f be a scalarly essentially \(\mu\) -integrable mapping of T into F such that, for every compact subset K of T, \(\mathbf{f}(\mathbf{K})\) is bounded. Then \(\int f d\mu\) belongs to F.

For, every bounded subset of F is relatively weakly compact (TVS, IV, §2, No.~2, Th. 1), and \(F = F''\) .

PROPOSITION 8. --- Let \(\mu\) be a bounded positive measure on T, S a \(\mu\) -measurable set carrying \(\mu\) , f a \(\mu\) -measurable mapping of T into F, such that f(S) is contained in a complete, bounded, balanced convex subset B of F. Then, f is scalarly \(\mu\) -integrable and \(\int f d\mu \in \mu(T)B \subset F\) .

Since S is \(\mu\) -integrable, there exists a partition of S formed by a \(\mu\) -negligible set N and a sequence \((\mathrm{K}_{n})\) of compact subsets such that the restriction of f to each \(K_{n}\) is continuous (Ch. IV, §4, No.~6, Cor. 3 of Th. 4 and §5, No.~1, Def. 1); \(\mathbf{f}(\mathrm{K}_n)\) is therefore a compact subset of F. The closed, balanced convex envelope \(\mathbf{B}_n\) of \(\mathbf{f}(\mathrm{K}_n)\) is then pre-compact (TVS, II, §4, No.~1, Prop. 3) and is contained in the complete subset B of F, therefore it is compact, and a fortiori weakly compact. Consequently (Cor. of Prop. 5) \(\mathbf{f}\varphi_{\mathrm{K}_n}\) is scalarly \(\mu\) -integrable, and

\[
\mathbf {z} _ {n} = \int \mathbf {f} \varphi_ {\mathrm{K} _ {n}} d \mu \in \mu (\mathrm{K} _ {n}) \mathrm{B} _ {n} \subset \mu (\mathrm{K} _ {n}) \mathrm{B}.
\]

For every continuous semi-norm p on F, it follows that

\[
p (\mathbf {z} _ {n}) \leqslant \mu (\mathrm{K} _ {n}) \cdot \sup _ {\mathbf {x} \in B} p (\mathbf {x});
\]

since B is bounded and since the series with general term \(\mu(\mathrm{K}_{n})\) is convergent and has sum \(\mu(\mathrm{T})\) , one sees that the sequence with general term \(s_{n}=z_{1}+z_{2}+\cdots+z_{n}\) is a Cauchy sequence in the complete subset \(\mu(\mathrm{T})\mathrm{B}\) of F. This sequence therefore converges to an element s of \(\mu(\mathrm{T})\mathrm{B}\) ; since one can suppose that \(\mathbf{f}(t)=0\) on T-S, Lebesgue's theorem applied to each of the functions \(\langle z', f\rangle\) ( \(z'\in F'\) ) proves that \(s=\int f d\mu\) .
